%-----------------------------------------------------------------------------------------------------------------------------------------------%
%	The MIT License (MIT)
%	Copyright (c) 2021 Jitin Nair
%-----------------------------------------------------------------------------------------------------------------------------------------------%

%----------------------------------------------------------------------------------------
%	DOCUMENT DEFINITION
%----------------------------------------------------------------------------------------

\documentclass[a4paper,12pt]{article}

%----------------------------------------------------------------------------------------
%	PACKAGES
%----------------------------------------------------------------------------------------
\usepackage{url}
\usepackage{parskip}

\RequirePackage{color}
\RequirePackage{graphicx}
\usepackage[usenames,dvipsnames]{xcolor}
\usepackage[scale=0.9]{geometry}
\usepackage{tabularx}
\usepackage{enumitem}

\newcolumntype{C}{>{\centering\arraybackslash}X}

\usepackage{supertabular}
\usepackage{tabularx}
\newlength{\fullcollw}
\setlength{\fullcollw}{0.47\textwidth}

\usepackage{titlesec}
\usepackage{multicol}
\usepackage{multirow}
\titleformat{\section}{\Large\scshape\raggedright}{}{0em}{}[\titlerule]
\titlespacing{\section}{0pt}{10pt}{10pt}

\usepackage[style=authoryear,sorting=ynt, maxbibnames=2]{biblatex}
\usepackage[unicode, draft=false]{hyperref}
\definecolor{linkcolour}{rgb}{0,0.2,0.6}
\hypersetup{colorlinks,breaklinks,urlcolor=linkcolour,linkcolor=linkcolour}

\usepackage{fontawesome5}

\newenvironment{jobshort}[2]
    {
    \begin{tabularx}{\linewidth}{@{}l X r@{}}
    \textbf{#1} & \hfill &  #2 \\[3.75pt]
    \end{tabularx}
    }
    {
    }

\newenvironment{joblong}[2]
    {
    \begin{tabularx}{\linewidth}{@{}l X r@{}}
    \textbf{#1} & \hfill &  #2 \\[3.75pt]
    \end{tabularx}
    \begin{minipage}[t]{\linewidth}
    \begin{itemize}[nosep,after=\strut, leftmargin=1em, itemsep=3pt,label=--]
    }
    {
    \end{itemize}
    \end{minipage}
    }

%----------------------------------------------------------------------------------------
%	BEGIN DOCUMENT
%----------------------------------------------------------------------------------------
\begin{document}

\pagestyle{empty}

%----------------------------------------------------------------------------------------
%	TITLE
%----------------------------------------------------------------------------------------
\begin{tabularx}{\linewidth}{@{} C @{}}
\Huge{刘飞月} \\[7.5pt]
\href{mailto:liufeiyue7@gmail.com}{\raisebox{-0.05\height}\faEnvelope \ liufeiyue7@gmail.com} \ $|$ \
\href{tel:+8617871637876}{\raisebox{-0.05\height}\faMobile \ +86 17871637876}
\end{tabularx}

%----------------------------------------------------------------------------------------
%	SUMMARY
%----------------------------------------------------------------------------------------
\section{Summary}
本人为安全科学与工程专业，主要研究方向为能源循环利用，聚焦热烟气相关工况，围绕烟气处理、能源回收过程中的安全防控、介质循环特性开展相关研究，兼顾能源高效利用与系统安全管控。

%----------------------------------------------------------------------------------------
%	EDUCATION
%----------------------------------------------------------------------------------------
\section{Education}
\begin{tabularx}{\linewidth}{@{}l X@{}}
2024 -- present & 硕士在读 安全科学与工程 at \textbf{太原理工大学} (211/双一流) \hfill \normalsize (GPA: 3.23/4.0) \\
2020 -- 2024 & 本科 安全工程 at \textbf{华北理工大学}
\end{tabularx}

%----------------------------------------------------------------------------------------
%	PROJECTS
%----------------------------------------------------------------------------------------
\section{Projects}
\begin{tabularx}{\linewidth}{ @{}l r@{} }
\textbf{深部高瓦斯难抽煤层热烟气强化增渗驱替机制及方法} & \hfill \\[3.75pt]
\multicolumn{2}{@{}X@{}}{国家自然科学基金重点项目。针对深部低渗高瓦斯煤层瓦斯抽采效率低、传统水力增透耗水量大的难题，以瓦斯电厂余热CO\textsubscript{2}-N\textsubscript{2}热烟气为循环介质，揭示热应力致裂、酸性矿物溶蚀协同增渗与CO\textsubscript{2}竞争吸附驱替瓦斯的多场耦合机理，建立煤体损伤--渗流--气体运移预测模型，研发低碳高效注气增透成套工艺，同步实现瓦斯灾害治理、烟气资源化循环与井下碳封存。} \\
\end{tabularx}

%----------------------------------------------------------------------------------------
%	SKILLS
%----------------------------------------------------------------------------------------
\section{Skills}
\begin{tabularx}{\linewidth}{@{}l X@{}}
实验设备 & \normalsize{工业分析仪、FT-IR 傅里叶变换红外光谱仪、NanoCT 扫描仪} \\
软件 & \normalsize{Avizo、Comsol、Origin} \\
数据分析 & \normalsize{CT 扫描数据、压汞数据、XRD 数据} \\
语言 & \normalsize{英语（大学英语六级证书）}
\end{tabularx}

\vfill
\center{\footnotesize Last updated: \today}

\end{document}
